% Validated Study Design staging fragment for ICNP 2026 venue draft.
% Paste or \input{} this fragment after:
% \section{Study Design}
% \label{sec:studydesign}
%
% This file contains only owner-validated reductions.

\subsection{Experimental Design}
\label{sec:expDesign}

Using the interface in \cref{sec:SystemModel}, this section defines the configuration axes used to answer \cref{sec:research_questions}. \Cref{tab:config_summary} summarizes the design dimensions, tested options, and RQ coverage.

\paragraph*{Network configuration}
We use the four-path diamond topology shown in \cref{fig:network_topology} and defined in \cref{subsec:topology}, with two-hop paths $P_1$ and $P_2$ and three-hop paths $P_3$ and $P_4$. This topology keeps the routing action space nontrivial while preserving tractable cross-product sweeps across \cref{tab:config_summary}.

Across all experiments, the total physical budget remains fixed at 35 qubits, matching the resource constraint in \cref{subsec:topology}; this prevents over-provisioning and creates the intended exploration--exploitation pressure. Per-path budgets are set by allocator policies in \cref{tab:setup-allocators}, enabling controlled comparison of static and adaptive resource management.

\begin{figure}[ht!]
\centering
\small
\begin{tikzpicture}[
  node distance=2.5cm and 3.0cm,
  thick,
  main/.style={circle,draw,minimum size=0.6cm,font=\small},
  link/.style={line width=0.8pt, draw=gray!55},
  path/.style={->,>=stealth,line width=1.1pt}
]
  % Nodes
  \node[main] (S) {$S$};
  \node[main, right=of S] (R1) {$R_1$};
  \node[main, below=of R1] (R2) {$R_2$};
  \node[main, right=of R1] (D) {$D$};

  % Physical links
  \draw[link] (S) -- (R1);
  \draw[link] (S) -- (R2);
  \draw[link] (R1) -- (D);
  \draw[link] (R2) -- (D);
  \draw[link] (R1) -- (R2);

  % P1: S -> R1 -> D (2-hop)
  \draw[path, blue]
    (S) to[bend left=15] node[above, font=\scriptsize] {$P_1$ (2-hop)} (R1)
        to[bend left=15] (D);

  % P2: S -> R2 -> D (2-hop)
  \draw[path, red]
    (S) to[bend right=15] node[left, font=\scriptsize] {$P_2$ (2-hop)} (R2)
        to[bend right=15] (D);

  % P3: S -> R1 -> R2 -> D (3-hop)
  \draw[path, green!60!black]
    (S) to[bend left=10] (R1)
        to[bend left=10] node[right, font=\scriptsize] {$P_3$ (3-hop)} (R2)
        to[bend left=10] (D);

  % P4: S -> R2 -> R1 -> D (3-hop)
  \draw[path, orange]
    (S) to[bend right=10] (R2)
        to[bend right=10] node[left, font=\scriptsize] {$P_4$ (3-hop)} (R1)
        to[bend right=10] (D);
\end{tikzpicture}
\caption{Four-path quantum-network testbed used for routing evaluation. Paths $P_1$ and $P_2$ have two hops, while $P_3$ and $P_4$ have three hops; per-path qubit budgets are varied by allocator policy.}
\label{fig:network_topology}
\end{figure}
